%! Project: Design and Conduct a Survey — Unit 1, Lesson 1.8 — Homework
% PILOT SHAPE (2026-09-06): modeled on AP Statistics lesson 1.4 minus its AP Practice page.
% This is the lesson's SCORED individual practice and the LAST component in the packet.
% Students START IT IN CLASS in the last ten minutes of the period, alone, while the
% teacher circulates for the second formative read; they finish it at home. Packet pages are
% the default; a Desmos activity may replace them on a given day (decided at Close & Assign).
% THIRD CONTEXT: the notes teach on the Student Council's survey, the Guided Practice runs on
% the class's own survey, and items 1-5 are on the Millbrook Public Library — a proper SRS of
% 80 against a front-desk clipboard. ITEM 6 IS THE PROJECT DELIVERABLE: the four-sentence
% report on the class's own data (it replaces the retired extensionbox).
% TWO PAGES MAX (user decision 2026-09-06): two boxes — items 1-2 on page 1, items 3-6 and
% the spiralbox on page 2 — so no item breaks across a page.
% Item 3 is the deliberate CONTRAST PAIR: the same scaling, once on the SRS (allowed) and once
% on the clipboard (not). Item 4 is the target misconception head-on (a clipboard that
% happens to match the SRS). Item 5 spirals to 1.4's response bias.
% THE DATA — Millbrook Public Library: 1,200 card holders; SRS of 80 by phone ->
%   28 book club (35%), 24 coding lab (30%), 20 story hour (25%), 8 job-search (10%);
%   0.35 x 1200 = 420; clipboard: 50 signed, 30 story hour -> 60%; 0.60 x 1200 = 720;
%   a clipboard at 13 of 50 = 26% (item 4).
% PAGE LOCKSTEP: the key is generated from this file; every work block is byte-identical.
\documentclass[12pt]{article}
\usepackage{saar-article}
\usepackage{saar-boxes}
\newcommand{\TallMath}[1]{$\displaystyle #1\rule[-1.4em]{0pt}{3.2em}$}
% Word bank strip for every fill-in cluster (user request 2026-09-06). Defined per document;
% the key inherits it and prints the same strip, so heights match.
\newcommand{\wordbank}[1]{\par\vspace{2pt}\noindent\fcolorbox{goldacc}{goldbg}{\parbox{\dimexpr\linewidth-2\fboxsep-2\fboxrule\relax}{\small\textbf{Word bank:} #1}}\par\vspace{4pt}}

\begin{document}

\pageheader{Unit 1, Lesson 1.8}{Homework}
% NAMESTRIP: no \namedateperiod here. The student writes their name once, on the
% cover this component is stapled behind.

\vspace{0.12in}

% Authored identically in homework/ and homework_key/.
\begin{remindbox}
\small
\textbf{This is your graded homework.} It is scored out of 12 and due at the start of the
first class after two study halls. You will start it in the last minutes of the period, so ask about anything that stalls
you before you leave, then finish it on your own. Item~6 is your project report.
\end{remindbox}

\vspace{0.08in}

\boxguard[20]
\begin{scenariobox}[Millbrook Public Library]{cerulean}
Millbrook has $1{,}200$ card holders and one grant to spend on a new program. The director draws
a \textbf{simple random sample of $80$} card holders from the membership list, reaches every one
by phone, and reads each the same item with four choices: evening book club $28$, teen coding
lab $24$, Saturday story hour $20$, job-search help $8$.

\vspace{2pt}
\textbf{Meanwhile}, a volunteer leaves a sign-up clipboard on the front desk for a week. $50$
people write something down: book club $10$, coding lab $6$, story hour $30$, job-search help $4$.
\end{scenariobox}

\vspace{0.08in}

\boxguard
\begin{notesbox}{Represent, Analyze, Report}
Every answer names the people it is about --- the $80$ phoned, the $50$ who signed, or all
$1{,}200$ card holders.
\wordbank{35\% \;\textbullet\; 30\% \;\textbullet\; 25\% \;\textbullet\; 10\% \;\textbullet\; 100\% \;\textbullet\; the share of the 80 \;\textbullet\; categorical}
\begin{enumerate}[leftmargin=1.6em, itemsep=6pt, topsep=2pt]

  \item Complete the relative frequency column for the phone survey.

        \begin{center}
        \renewcommand{\arraystretch}{1.12}
        \begin{tabularx}{\linewidth}{@{} X c p{5.4cm} @{}}
        \toprule
        \textbf{Program} & \textbf{Frequency} & \textbf{Relative frequency} \\
        \midrule
        Evening book club   & $28$ & \blank{4.4cm} \\
        Teen coding lab     & $24$ & \blank{4.4cm} \\
        Saturday story hour & $20$ & \blank{4.4cm} \\
        Job-search help     & $8$  & \blank{4.4cm} \\
        \midrule
        \textbf{Total}      & $80$ & \blank{4.4cm} \\
        \bottomrule
        \end{tabularx}
        \end{center}

  \item The director's bar graph of the phone survey.

        \begin{center}
        \begin{tikzpicture}[x=0.95cm, y=0.046cm]
          \foreach \y in {10,20,30,40}
            \draw[linegray] (0,\y) -- (10.4,\y);
          \foreach \y in {0,10,20,30,40}
            \node[left, font=\scriptsize] at (-0.05,\y) {\y\%};
          \foreach \i/\h in {1/35, 2/30, 3/25, 4/10}
            \fill[ceruleanlight] ({2.6*\i-2.0},0) rectangle ({2.6*\i-0.6},\h);
          \foreach \i/\lab in {1/Book club, 2/Coding lab, 3/Story hour, 4/Job-search}
            \node[below, font=\scriptsize, align=center, text width=2.4cm]
              at ({2.6*\i-1.3},-0.5) {\lab};
          \draw[cerulean, thick] (0,42) -- (0,0) -- (10.4,0);
        \end{tikzpicture}
        \end{center}

        The height of the first bar tells a reader \blank{4.4cm}.
        \par\vspace{6pt}
        The bars are separated because ``program'' is a \blank{3.4cm} variable.

\end{enumerate}
\end{notesbox}

\boxguard[30]
\begin{notesbox}{Represent, Analyze, Report, continued}
\begin{enumerate}[leftmargin=1.6em, itemsep=3pt, topsep=0pt, start=3]

  \item Two surveys, one library, the same arithmetic.

        \textbf{(a)} Scale the phone survey's book-club percent to all $1{,}200$ card holders.

        \begin{work}
          0.35 \times 1200 &= 420 \text{ card holders}
        \end{work}

        \textbf{(b)} What percent of the $50$ who signed the clipboard chose story hour?

        \begin{work}
          \dfrac{30}{50} &= 0.60 = 60\%
        \end{work}

        \textbf{(c)} One of these percents may be scaled to all $1{,}200$; one may not. Which, and why?
        \par\writelines{2}

  \item Suppose the clipboard had shown $13$ of $50$ for story hour --- $26\%$, almost exactly the
        phone survey's $25\%$. Does that near-match make it trustworthy for all $1{,}200$? Explain.
        \par\writelines{2}

  \item An early draft of the director's item read: \emph{``Wouldn't you love a coding lab for
        our teens?''} Name the instrument rule it breaks and the name Lesson 1.4 gave that
        bias, then repair it.
        \par\writelines{2}

  \item \textbf{Your project report.} From your class's tally in the notes, write the four
        sentences: what you asked and of whom; how the sample was chosen and how many answered;
        what you found (the two largest categories, as percents); what it means --- \emph{about
        the people your design actually covers}.
        \par\writelines{4}

\end{enumerate}
\end{notesbox}

\boxguard[3]   % the three-line spiralbox just fits the foot of page 2; a larger guard pushes it to a third page
\begin{spiralbox}
\textbf{Unit 1 is finished.} Every idea in it answered one question: how much should you believe
a number somebody hands you? Next: review, the unit test, then Unit 2 --- one \emph{quantitative}
variable, where the bar graph gives way to a dot plot and the bars start to touch.
\end{spiralbox}

\end{document}
